% GENERATED FILE. Edit canonical lessons, then run npm run generate:books.
% canonical-source-hash: fnv1a64:bfac03588c1023ae
% canonical-lessons: ES-C692-sembrar, ES-C692-regar, ES-C692-cosechar, ES-C692-pozo, ES-C692-florecer

\chapter{Sembrar y regar}
\label{ch:es-sembrar-y-regar}
\hlchaptermodality{692}

\begin{chapteropening}
\textbf{By the end of this chapter:} \emph{I can describe farm work: sowing, watering, harvesting, the well and the first flowers.}

Complete the last lesson of chapter 692: I can describe farm work: sowing, watering, harvesting, the well and the first flowers.
\end{chapteropening}

\section[sembrar]{sembrar — to sow}

\label{lesson:ES-C692-sembrar}



\begin{quote}
Before the new one: say the Spanish for the vineyard, then the Spanish for to grow.
\end{quote}

\subsection*{You'll want to know: sembrar}
\textbf{sembrar} — \textquotedblleft{}to sow\textquotedblright{}.

\begin{quote}
En primavera siembran el maíz.
\end{quote}

\begin{quote}
\emph{In spring they sow the maize.}
\end{quote}

The e becomes ie under stress: \emph{siembro}, \emph{siembra}.

\begin{cousinweb}
\emph{sembrar} is inherited from Latin \emph{sēmināre} \textquotedblleft{}to sow\textquotedblright{} (EtymDB), from \emph{sēmen} \textquotedblleft{}seed\textquotedblright{}.
\end{cousinweb}

\begin{grammarlens}[title={What you've built}]
One more word for telling, explaining and describing at B1.
\end{grammarlens}

\subsection*{Your turn}
\begin{itemize}
\raggedright
  \item \emph{Say it:} \emph{sembrar}
  \item \emph{Say it:} \emph{sembrar}, once more
  \item \emph{Say it:} a sentence of your own with \emph{sembrar}
  \item \emph{Recall:} say the Spanish for the vineyard, then the Spanish for to grow, then say \emph{sembrar} again
\end{itemize}

\begin{tcolorbox}[breakable,colback=teal!4,colframe=teal!35!black,title={Before you move on}]
What does sembrar mean? (\textquotedblleft{}To sow\textquotedblright{}.) Say it once more.
\end{tcolorbox}
\section[regar]{regar — to water (plants)}

\label{lesson:ES-C692-regar}



\begin{quote}
Before the new one: say the Spanish for to grow, then the Spanish for to sow.
\end{quote}

\subsection*{You'll want to know: regar}
\textbf{regar} — \textquotedblleft{}to water (plants)\textquotedblright{}.

\begin{quote}
Riego las plantas por la noche.
\end{quote}

\begin{quote}
\emph{I water the plants in the evening.}
\end{quote}

The e becomes ie under stress: \emph{riego}, \emph{riega}.

\begin{cousinweb}
\emph{regar} is inherited from Latin \emph{rigāre} \textquotedblleft{}to water, to moisten\textquotedblright{} (EtymDB).
\end{cousinweb}

\begin{grammarlens}[title={What you've built}]
One more word for telling, explaining and describing at B1.
\end{grammarlens}

\subsection*{Your turn}
\begin{itemize}
\raggedright
  \item \emph{Say it:} \emph{regar}
  \item \emph{Say it:} \emph{regar}, once more
  \item \emph{Say it:} a sentence of your own with \emph{regar}
  \item \emph{Recall:} say the Spanish for to grow, then the Spanish for to sow, then say \emph{regar} again
\end{itemize}

\begin{tcolorbox}[breakable,colback=teal!4,colframe=teal!35!black,title={Before you move on}]
What does regar mean? (\textquotedblleft{}To water (plants)\textquotedblright{}.) Say it once more.
\end{tcolorbox}
\section[cosechar]{cosechar — to harvest, to reap}

\label{lesson:ES-C692-cosechar}



\begin{quote}
Before the new one: say the Spanish for to sow, then the Spanish for to water.
\end{quote}

\subsection*{You'll want to know: cosechar}
\textbf{cosechar} — \textquotedblleft{}to harvest, to reap\textquotedblright{}.

\begin{quote}
En septiembre cosechan la uva.
\end{quote}

\begin{quote}
\emph{In September they harvest the grapes.}
\end{quote}

\emph{La cosecha} is the harvest.

\begin{cousinweb}
\emph{cosechar} is built on \emph{cosecha} \textquotedblleft{}harvest\textquotedblright{}, which comes from Latin \emph{collēcta} \textquotedblleft{}a gathering\textquotedblright{} (EtymDB), from \emph{colligere} \textquotedblleft{}to gather\textquotedblright{}.
\end{cousinweb}

\begin{grammarlens}[title={What you've built}]
One more word for telling, explaining and describing at B1.
\end{grammarlens}

\subsection*{Your turn}
\begin{itemize}
\raggedright
  \item \emph{Say it:} \emph{cosechar}
  \item \emph{Say it:} \emph{cosechar}, once more
  \item \emph{Say it:} a sentence of your own with \emph{cosechar}
  \item \emph{Recall:} say the Spanish for to sow, then the Spanish for to water, then say \emph{cosechar} again
\end{itemize}

\begin{tcolorbox}[breakable,colback=teal!4,colframe=teal!35!black,title={Before you move on}]
What does cosechar mean? (\textquotedblleft{}To harvest, to reap\textquotedblright{}.) Say it once more.
\end{tcolorbox}
\section[el pozo]{el pozo — the well}

\label{lesson:ES-C692-pozo}



\begin{quote}
Before the new one: say the Spanish for to water, then the Spanish for to harvest.
\end{quote}

\subsection*{You'll want to know: el pozo}
\textbf{el pozo} — \textquotedblleft{}the well\textquotedblright{}.

\begin{quote}
Sacan agua del pozo.
\end{quote}

\begin{quote}
\emph{They draw water from the well.}
\end{quote}

\emph{Un pozo de petróleo} — an oil well.

\begin{cousinweb}
\emph{pozo} is inherited from Latin \emph{puteus} \textquotedblleft{}a well, a pit\textquotedblright{} (EtymDB).
\end{cousinweb}

\begin{grammarlens}[title={What you've built}]
One more word for telling, explaining and describing at B1.
\end{grammarlens}

\subsection*{Your turn}
\begin{itemize}
\raggedright
  \item \emph{Say it:} \emph{el pozo}
  \item \emph{Say it:} \emph{el pozo}, once more
  \item \emph{Say it:} a sentence of your own with \emph{el pozo}
  \item \emph{Recall:} say the Spanish for to water, then the Spanish for to harvest, then say \emph{el pozo} again
\end{itemize}

\begin{tcolorbox}[breakable,colback=teal!4,colframe=teal!35!black,title={Before you move on}]
What does el pozo mean? (\textquotedblleft{}The well\textquotedblright{}.) Say it once more.
\end{tcolorbox}
\section[florecer]{florecer — to flower, to bloom; to flourish}

\label{lesson:ES-C692-florecer}



\begin{quote}
Before the new one: say the Spanish for to harvest, then the Spanish for the well.
\end{quote}

\subsection*{You'll want to know: florecer}
\textbf{florecer} — \textquotedblleft{}to flower, to bloom; to flourish\textquotedblright{}.

\begin{quote}
Los árboles florecen en primavera.
\end{quote}

\begin{quote}
\emph{The trees blossom in spring.}
\end{quote}

It adds a z: \emph{florezco}, \emph{florezca}.

\begin{cousinweb}
\emph{florecer} is inherited from Latin \emph{flōrēscere} \textquotedblleft{}to begin to bloom\textquotedblright{} (EtymDB), from \emph{flōs, flōris} \textquotedblleft{}flower\textquotedblright{}.
\end{cousinweb}

\begin{grammarlens}[title={What you've built}]
One more word for telling, explaining and describing at B1.
\end{grammarlens}

\subsection*{Your turn}
\begin{itemize}
\raggedright
  \item \emph{Say it:} \emph{florecer}
  \item \emph{Say it:} \emph{florecer}, once more
  \item \emph{Say it:} a sentence of your own with \emph{florecer}
  \item \emph{Recall:} say the Spanish for to harvest, then the Spanish for the well, then say \emph{florecer} again
\end{itemize}

\begin{tcolorbox}[breakable,colback=teal!4,colframe=teal!35!black,title={Before you move on}]
What does florecer mean? (\textquotedblleft{}To flower, to bloom; to flourish\textquotedblright{}.) Say it once more.
\end{tcolorbox}
